\documentclass[12pt]{article}

\usepackage[margin=1in]{geometry}
\usepackage{amsmath,amssymb}
\usepackage{microtype}
\usepackage{parskip}

\title{What Would It Mean to Differentiate Marching Cubes?}
\author{}
\date{}

\begin{document}

\maketitle

\section*{1. Marching cubes as a function}

For fixed process parameters
$\boldsymbol{\theta}=(P,\sigma,v)^{\mathsf T}$, the heat-transfer model
provides temperatures on a three-dimensional grid,
\begin{equation}
  T_{ijk}=T(x_i,y_j,z_k;\boldsymbol{\theta}).
\end{equation}
Marching cubes may be viewed abstractly as the function
\begin{equation}
  \mathcal M:
  \left(
    \{T_{ijk}\},\,T_{\mathrm{liq}},\,\text{grid geometry}
  \right)
  \longmapsto
  (V,F),
  \label{eq:marching-cubes-map}
\end{equation}
where $V$ contains the vertices of the extracted liquidus isosurface and
$F$ contains its triangular faces. The process parameters are not explicit
inputs to $\mathcal M$. Their effects enter through the sampled temperatures
$T_{ijk}(\boldsymbol{\theta})$.

The melt-pool dimensions are subsequently calculated from the isosurface
vertices. If $V_y$ and $V_z$ denote their transverse and vertical
coordinates, respectively, then
\begin{equation}
  w=\frac{\max(V_y)-\min(V_y)}{2},
  \qquad
  d=-\min(V_z).
  \label{eq:mesh-dimensions}
\end{equation}
Thus, the complete numerical measurement may be represented as
\begin{equation}
  \boldsymbol{\theta}
  \longrightarrow
  \{T_{ijk}\}
  \longrightarrow
  \mathcal M(\{T_{ijk}\})
  \longrightarrow
  (w,d).
  \label{eq:measurement-chain}
\end{equation}

\section*{2. What direct AD through marching cubes would require}

In an OTI calculation, every input temperature would carry both its real
value and its process-parameter derivatives. Schematically,
\begin{equation}
  T_{ijk}^{\mathrm{OTI}}
  =
  T_{ijk}
  +\varepsilon_P\frac{\partial T_{ijk}}{\partial P}
  +\varepsilon_\sigma\frac{\partial T_{ijk}}{\partial\sigma}
  +\varepsilon_v\frac{\partial T_{ijk}}{\partial v}.
  \label{eq:oti-temperature}
\end{equation}

Marching cubes first compares the temperatures at the corners of each grid
cell with $T_{\mathrm{liq}}$ to determine which cell edges intersect the
liquidus surface. For an intersected edge with endpoints $a$ and $b$, a
surface vertex is obtained by interpolation:
\begin{align}
  \lambda
  &=\frac{T_{\mathrm{liq}}-T_a}{T_b-T_a},\\
  \mathbf{x}_{\mathrm{surf}}
  &=\mathbf{x}_a
    +\lambda(\mathbf{x}_b-\mathbf{x}_a).
  \label{eq:edge-interpolation}
\end{align}
If $T_a$ and $T_b$ are OTI quantities and the arithmetic in
Equation~\eqref{eq:edge-interpolation} is performed using OTI operations,
then the resulting vertex coordinates also carry derivatives:
\begin{equation}
  \mathbf{x}_{\mathrm{surf}}^{\mathrm{OTI}}
  =
  \mathbf{x}_{\mathrm{surf}}
  +\varepsilon_P
    \frac{\partial\mathbf{x}_{\mathrm{surf}}}{\partial P}
  +\varepsilon_\sigma
    \frac{\partial\mathbf{x}_{\mathrm{surf}}}{\partial\sigma}
  +\varepsilon_v
    \frac{\partial\mathbf{x}_{\mathrm{surf}}}{\partial v}.
\end{equation}

Obtaining $\partial w/\partial\theta_j$ and
$\partial d/\partial\theta_j$ directly would then require propagating these
OTI-valued vertices through the minimum and maximum operations in
Equation~\eqref{eq:mesh-dimensions}. Therefore, ``AD of marching cubes''
would actually mean differentiating the entire chain
\begin{equation}
  \boxed{
  \text{OTI temperature grid}
  \longrightarrow
  \text{isosurface vertices}
  \longrightarrow
  \text{extremal vertices}
  \longrightarrow
  (w,d)}.
\end{equation}

\section*{3. Why direct differentiation is awkward}

The interpolation in Equation~\eqref{eq:edge-interpolation} is differentiable
as long as the same edge continues to contain the surface intersection. Other
parts of marching cubes are discrete:
\begin{itemize}
  \item comparisons with $T_{\mathrm{liq}}$ select one of the possible cell
        configurations;
  \item the selected configuration determines the number and connectivity of
        the surface triangles; and
  \item the minimum and maximum operations select which mesh vertices define
        the reported dimensions.
\end{itemize}
AD can propagate derivatives through one fixed selection of these branches.
However, the derivative can be discontinuous or undefined when a perturbation
changes the marching-cubes configuration or causes a different vertex to
become extremal. A standard marching-cubes implementation also expects an
ordinary real-valued array and therefore cannot directly operate on OTI
objects without being rewritten or adapted.

\section*{4. The implicit alternative}

The present method avoids differentiating this discrete geometry pipeline.
Define
\begin{equation}
  F(\mathbf{x};\boldsymbol{\theta})
  =T(\mathbf{x};\boldsymbol{\theta})-T_{\mathrm{liq}}.
\end{equation}
Every point $\mathbf{x}^*(\boldsymbol{\theta})$ on the liquidus surface
satisfies
\begin{equation}
  F(\mathbf{x}^*(\boldsymbol{\theta});\boldsymbol{\theta})=0.
\end{equation}
Differentiation with respect to a process parameter $\theta_j$ gives
\begin{equation}
  \nabla_{\mathbf{x}}F\mathbin{\cdot}
  \frac{\partial\mathbf{x}^*}{\partial\theta_j}
  +F_{\theta_j}=0.
  \label{eq:implicit-general}
\end{equation}

At a transverse extremum, the surface normal is aligned with the
$y$-direction. Equation~\eqref{eq:implicit-general} therefore reduces to
\begin{equation}
  \frac{\partial y^*}{\partial\theta_j}
  =-\frac{T_{\theta_j}}{T_y}.
  \label{eq:transverse-sensitivity}
\end{equation}
At the deepest point,
\begin{equation}
  \frac{\partial z_{\min}}{\partial\theta_j}
  =-\frac{T_{\theta_j}}{T_z},
  \qquad
  \frac{\partial d}{\partial\theta_j}
  =\frac{T_{\theta_j}}{T_z},
  \label{eq:depth-sensitivity}
\end{equation}
because $d=-z_{\min}$. Combining the two transverse extrema gives
\begin{equation}
  \frac{\partial w}{\partial\theta_j}
  =\frac{1}{2}
  \left[
    -\left.\frac{T_{\theta_j}}{T_y}\right|_{\mathbf{x}_{y^+}}
    +\left.\frac{T_{\theta_j}}{T_y}\right|_{\mathbf{x}_{y^-}}
  \right].
  \label{eq:width-sensitivity}
\end{equation}

The ratio in Equations~\eqref{eq:transverse-sensitivity}--
\eqref{eq:width-sensitivity} has a simple interpretation:
\begin{equation}
  \text{boundary displacement}
  =-\frac{\text{temperature change caused by the parameter}}
          {\text{temperature change per unit distance}}.
\end{equation}
The OTI parameter derivative supplies the numerator, while the spatial
temperature gradient converts that temperature change into movement of the
liquidus boundary.

\section*{5. Original and current implementations}

In the original implementation, the complete local grid was evaluated using
OTI arithmetic. Marching cubes was applied only to the real temperature field,
and the OTI derivative fields were interpolated to the three extremal points.
The interpolated derivatives were then converted into geometric sensitivities
using Equations~\eqref{eq:depth-sensitivity} and
\eqref{eq:width-sensitivity}.

The current implementation first evaluates the complete real-valued grid to
locate the liquidus extrema. It then evaluates OTI quantities only on the
small interpolation stencils surrounding those points. This preserves the
same local sensitivity calculation without differentiating marching cubes or
carrying OTI quantities over the entire bounding box.

\end{document}
